\chapter{Recopilación de Información}
\label{ap:recopilacion-informacion}

Este apéndice reúne los antecedentes utilizados para caracterizar el trabajo de
clasificación acústica, formular los requerimientos de F.A.M.A. y documentar
decisiones durante su desarrollo. La información procede de la revisión de
documentos técnicos y del análisis de registros de experimentación, pruebas y
operación del sistema.

\section{Revisión documental}

La revisión de antecedentes permitió describir el flujo de trabajo del
investigador, desde la obtención de grabaciones etiquetadas hasta la
preparación de los datos, el entrenamiento y la revisión de predicciones.
En la descripción del problema se identificaron dificultades para conservar
la procedencia de los audios y relacionar las ejecuciones con sus resultados.
Estos antecedentes sirvieron para delimitar las funciones de ingesta,
preprocesamiento, entrenamiento, predicción y retroalimentación descritas en
el capítulo~3.

Se consultó la documentación de Google Cloud Storage, Cloud SQL y gestión de
identidades y accesos, recogida en las referencias de la memoria. La revisión
permitió identificar qué recursos se necesitan para conservar los archivos
de audio, registrar sus metadatos y autorizar al nodo de procesamiento.
También se revisaron los antecedentes técnicos sobre representación de audio
y aprendizaje automático que sustentan la preparación de las señales para
el entrenamiento local.

La documentación interna complementó estas fuentes con decisiones propias
del proyecto. El registro ADR~0013 describe la captura de confirmaciones y
correcciones y la revisión de aportes antes de su incorporación. El
ADR~0069 registra el acuerdo de mantener conjuntos separados de
entrenamiento, validación y prueba, conservar juntos los grupos de
grabaciones relacionados y reservar la activación del modelo predeterminado
para una acción administrativa.

\section{Registros de experimentación y observación técnica}

En la prueba de concepto bioacústica, el equipo examinó grabaciones
etiquetadas procedentes de Xeno-canto. El informe de septiembre de 2026
conserva los criterios de selección de especies, calidad de los audios
y procedencia de las grabaciones. El análisis del origen de los
ejemplos permitió identificar un riesgo para la evaluación cuando audios
relacionados quedan repartidos entre el entrenamiento y la prueba. Las
decisiones posteriores fijaron la selección y distribución de grabaciones
mediante versiones de particiones.

El registro de integración del 24 de octubre de 2026 reúne comprobaciones
realizadas por el equipo en un entorno de laboratorio, con datos sintéticos
y bases PostgreSQL temporales. Se registraron resultados relativos a
permisos, particiones, persistencia y activación de modelos. Estos ensayos
aportaron información sobre el comportamiento del software bajo condiciones
controladas.

El registro de mantenimiento del 25 de octubre de 2026 documenta
comprobaciones en el nodo local del proyecto conectado a los servicios
oficiales. Se revisaron las conexiones con PostgreSQL y Google Cloud Storage,
el acceso administrativo, el catálogo de modelos y el esquema de base de
datos. También se contrastaron las copias restauradas en bases locales
separadas. Estas observaciones sirvieron para detallar los pasos de
preparación, verificación y recuperación del capítulo~8.

\section{Aplicación de la información recopilada}

\begin{description}
    \item[Gestión de datos] La descripción del trabajo con grabaciones y la
    revisión de los servicios cloud sustentaron la separación entre el
    almacenamiento de archivos, el registro de metadatos y el procesamiento
    local. Esta decisión se refleja en los requerimientos \texttt{RF\_01} y
    \texttt{RF\_02}.

    \item[Trazabilidad experimental] La revisión de los datos y las pruebas
    de preparación respaldaron el registro de las particiones y su relación
    con los modelos entrenados. Esta información se utiliza en el recorrido
    de entrenamiento y evaluación asociado a \texttt{RF\_04}.

    \item[Revisión de predicciones] La necesidad de registrar la
    confirmación o corrección de una clase se formalizó en
    \texttt{RF\_06}. El registro de decisiones distingue la captura de
    retroalimentación, la revisión del aporte y su uso en una preparación
    posterior.

    \item[Operación del sistema] Los registros de acceso, conexión y
    recuperación permitieron precisar las tareas del administrador y del
    responsable técnico del nodo. La verificación del laboratorio y la
    comprobación en el nodo local se documentaron por separado, junto con
    las condiciones previstas para la implantación institucional.
\end{description}

\section{Respaldo documental y alcance}

Los antecedentes se conservan en el repositorio del proyecto. El informe
\emph{Primera prueba de concepto} respalda la selección de grabaciones.
Los ADR~0013 y 0069 registran las decisiones sobre retroalimentación y
preparación experimental. Los registros de integración del 24 de octubre y
de actualización del nodo oficial del 25 de octubre documentan el entorno
de comprobación, las operaciones observadas y sus resultados.

El plan de trabajo del objetivo específico~2 contempla entrevistas y
reuniones de avance con el profesor guía y con investigadores del área.
Su ejecución queda por acreditar mediante los instrumentos, actas o
registros correspondientes. La evidencia reunida en este apéndice procede
de las fuentes documentales y de los registros técnicos descritos.

% Rutas de respaldo en el repositorio de F.A.M.A.
% docs/modelos/01_primera_prueba_poc.md
% docs/adr/0013_bucle_retroalimentacion_curacion_semi_manual_mlops.md
% docs/adr/0069-integracion-de-acceso-evaluacion-y-datasets.md
% docs/manuales/particiones_versionadas.md
% docs/manuales/integracion_dev_20261024.md
% docs/receipts/integracion_dev_20261024.json
% docs/manuales/actualizacion_oficial_0004_20261025.md
% docs/receipts/actualizacion_oficial_0004_20261025.json
